\section{Delay functions and the continuous path process}
\label{sec:functions}

We first analyze the continuous model in which a prime exponent is chosen
with harmonic measure \(dx/x\).  The linear-sieve functions provide its
benchmark, and their derivatives will convert harmonic path weights into
probability weights.  We then use regeneration to evaluate compact-endpoint
sums at scale \(B^{-2}\) and to control their tails.  These are the
estimates needed to propagate the reference tree's boundary discrepancy.

\subsection{Future integrals for the sieve functions}

We use the classical lower and upper linear-sieve functions; see
\cite[Section~21.2.2]{MontgomeryVaughanII}. The future-integral and path
representations needed below are developed here.

Write \(A=2e^\gamma\).  Define the continuous delay functions \(f,F\) by
\[
 f(s)=0\quad(0\le s\le2),\qquad sF(s)=A\quad(1\le s\le3),
\]
and
\[
 (sf(s))'=F(s-1)\quad(s>2),\qquad
 (sF(s))'=f(s-1)\quad(s>3).
\]
These prescriptions determine the functions successively on intervals of
length one.  We shall also use \(F(s)=A/s\) on \(0.95\le s<1\).
Only at a starting lower state, and when defining its derivative weight,
we continue \(f\) below \(2\) by
\begin{equation}\label{eq:functions-starting-extension}
 f_{\rm ext}(s)=\frac{A\log(s-1)}s\qquad(1.98\le s<2).
\end{equation}
This is the same expression as the ordinary \(f\) on \(2\le s\le4\).
An omitted lower child of ratio less than \(2\) always has benchmark
\(f=0\); the continuation \eqref{eq:functions-starting-extension} is not
used at such a child.

Put \(f_\ee=f\) and \(f_\oo=F\).  The continuous benchmark at a state
of type \(i\), gap \(r\), and cutoff \(b\) is \(f_i(r/b)/b\).
For fixed \(r\), the delay equations give
\begin{equation}\label{eq:functions-benchmark-cutoff}
 \frac{d}{db}\left(\frac{f_i(r/b)}b\right)
       =-\frac{f_{i'}(r/b-1)}{b^2},
\end{equation}
where \(i'\) is the opposite type.  On the right, the value at an
omitted lower child is defined to be zero, including when its formal
ratio is negative.  On the ordinary ranges the identity holds for
\(r/b>2\) at a lower state and \(r/b>1\) at an upper state.
At an upper state the right side is zero when \(r/b-1\le2\), exactly
as required for an omitted lower child.  Integrating with respect to the
cutoff therefore gives the standard continuous tree recurrence above
exponent \(2\), with the boundary value at cutoff \(2\) left to be
specified.  The explicit starting extensions satisfy the same identity
in the interiors of their ranges, with one-sided derivatives at the
endpoints; the extension of \(f\) is used at the parent lower state,
never at an omitted child.

The root lies at \(r_0/B=2-a_\star/B+o(B^{-1})\), where the ordinary
lower-sieve function vanishes.  Its continued benchmark is
\[
 \frac{f_{\rm ext}(r_0/B)}B
       =-a_\star e^\gamma B^{-2}+o(B^{-2}),
\]
since \(f'_{\rm ext}(2)=e^\gamma\).  Thus a positive reference margin
must come from the finite boundary, after its discrepancy has been
propagated through all larger exponents.  To evaluate that contribution,
we will use the positive derivative weights of \(f,F\) to normalize the
harmonic transitions.  The first step is to establish positivity and
decay in a form that justifies this change of measure.

Let \(\rho\) be the Dickman function, determined by
\[
 \rho(u)=1\quad(0\le u\le1),\qquad
 u\rho'(u)=-\rho(u-1)\quad(u>1).
\]
For the unextended functions put
\[
 a(s)=1-f(s)\quad(s\ge0),\qquad
 c(s)=F(s)-1\quad(s\ge1),\qquad
 c_0(s)=\frac{A\rho(s-1)}s\quad(s\ge1).
\]

\begin{lemma}[Future integrals]\label{lem:future-integrals}
The functions \(a,c\) are positive and satisfy, with absolute comparison
constants,
\begin{align}
 a(s)+c(s)&=c_0(s),&
 a(s)&\asymp c_0(s),&
 c(s)&\asymp c_0(s)\qquad(s\ge1),\label{eq:functions-deficits}\\
 s a(s)&=\int_{s-1}^{\infty}c(t)\,dt
       &&(s\ge2),\label{eq:functions-future-a}\\
 s c(s)&=\int_{s-1}^{\infty}a(t)\,dt
       &&(s\ge1).\label{eq:functions-future-c}
\end{align}
For \(u\ge2\),
\begin{equation}\label{eq:functions-dickman-ratio}
 \frac1{4u^2}\le\frac{\rho(u)}{\rho(u-1)}\le\frac1u .
\end{equation}
Moreover, for each fixed \(\lambda>0\), there is \(U_\lambda\) such that
\begin{equation}\label{eq:functions-dickman-logderivative}
 -\frac{\rho'(u)}{\rho(u)}\ge\lambda\qquad(u\ge U_\lambda).
\end{equation}
\end{lemma}

\begin{proof}
The delay equation gives the integral identity
\begin{equation}\label{eq:functions-dickman-integral}
 u\rho(u)=\int_{u-1}^{u}\rho(v)\,dv\qquad(u\ge1).
\end{equation}
For example, both sides agree at \(u=1\), and their derivatives agree
where the delay equation applies.  The identity and the delay equation
give positivity and monotonicity: a first zero of \(\rho\) would contradict
the positive integral of its immediately preceding values.
Monotonicity gives the upper bound in
\eqref{eq:functions-dickman-ratio}.  Integrating first over the lower
half of each of two unit intervals gives
\[
 \rho(u)\ge\frac{\rho(u-\tfrac12)}{2u}
 \ge\frac{\rho(u-1)}{4u(u-\tfrac12)}
 \ge\frac{\rho(u-1)}{4u^2}.
\]
The upper ratio bound also shows that \(\rho\) tends to zero faster
than any exponential.

The negative logarithmic derivative is at least \(1\).  If it is at
least \(\lambda\) on a tail, then, on a later tail,
\[
 \int_{u-1}^{u}\rho(v)\,dv
 \le \rho(u-1)\frac{1-e^{-\lambda}}{\lambda}.
\]
Equation \eqref{eq:functions-dickman-integral} consequently improves
the lower bound to \(\lambda/(1-e^{-\lambda})\).
Iteration tends to infinity, since this map is strictly larger than its
argument and has no finite positive fixed point.  This proves
\eqref{eq:functions-dickman-logderivative}.

Directly from the delay equations and the initial conditions,
\begin{equation}\label{eq:functions-gap-sum}
 F(s)-f(s)=c_0(s),\qquad
 F(s)+f(s)=A\omega_{\rm Bu}(s)\qquad(s\ge1).
\end{equation}
Here the continuous Buchstab function is determined by
\(s\omega_{\rm Bu}(s)=1\) on \([1,2]\) and
\((s\omega_{\rm Bu}(s))'=\omega_{\rm Bu}(s-1)\) above \(2\);
its limit is stated in Section~\ref{sec:inputs}.
For the first equality, the equation for
\(s(F-f)\) is the negative delay equation for \(A\rho(s-1)\);
for the second, it is the Buchstab equation with initial value \(A/s\).
We next prove the individual positivity and future integrals; these do
not follow merely by taking the difference in
\eqref{eq:functions-gap-sum}.

To prove these properties, we use an auxiliary absorbing chain.  This
chain establishes positivity of the deficits; the harmonic-path chain
will instead be defined from the derivative weights below.
For \(s\ge2\), give the auxiliary chain transition density
\begin{equation}\label{eq:functions-auxiliary-kernel}
 \frac{c_0(t)}{s c_0(s)}\,\1_{\{t\ge s-1\}}\,dt,
\end{equation}
and absorb it on reaching \([1,2)\).  Its normalization is
\[
 \int_{s-1}^{\infty}c_0(t)\,dt
 =A\int_{s-2}^{\infty}\frac{\rho(v)}{v+1}\,dv
 =A\rho(s-1)=s c_0(s).
\]
The integration at infinity is justified by the decay already proved.
The excess of the next state above \(s-1\) has tail
\[
 \frac{\rho(s-1+h)}{\rho(s-1)}\qquad(h\ge0).
\]
For large \(s\), \eqref{eq:functions-dickman-logderivative} bounds this
tail by \(e^{-\lambda h}\), with \(\lambda\) as large as desired.
For \(\lambda>1\), therefore,
\[
 \E(e^{s_{\rm next}}\mid s)
 \le e^{s-1}\frac{\lambda}{\lambda-1}.
\]
Taking \(\lambda=2\) makes this a strict contraction of \(V(s)=e^s\),
with factor \(2/e<1\), outside a fixed compact.
Stopped drift implies almost sure return to that compact.  From any
state in the compact, successive draws in
\([s-1,s-1+1/4]\) reach \([1,2)\) within a bounded number of steps,
with probability bounded below independently of the initial state
in the compact.  The densities in these intervals have a positive
compact lower bound.  Repeating the attempt after each return proves
almost sure absorption.

Alternate labels \(a,c\) at every draw.  At the terminal state \(t\)
assign fractions \(1/c_0(t)\) and \(1-1/c_0(t)\) to the two labels,
respectively.  Since \(c_0(t)=A/t\) on \([1,2)\), both terminal
fractions lie in a fixed compact subinterval of \((0,1)\).
The expected terminal fraction, multiplied by \(c_0(s)\), defines
positive functions \(a^\circ(s),c^\circ(s)\) with
\[
 a^\circ+c^\circ=c_0,\qquad
 a^\circ,c^\circ\asymp c_0,
\]
and terminal values \(a^\circ=1,\ c^\circ=c_0-1\) on \([1,2)\).
The first-step equations give
\[
 s a^\circ(s)=\int_{s-1}^{\infty}c^\circ(t)\,dt,\qquad
 s c^\circ(s)=\int_{s-1}^{\infty}a^\circ(t)\,dt
 \qquad(s\ge2).
\]
In particular these functions are continuous from the right at \(2\).

Set \(U=(2+c^\circ-a^\circ)/A\).  The last two equations imply
\((sU(s))'=U(s-1)\) above \(2\), in integrated form, while
\(sU(s)=1\) on \([1,2)\).
Also \(U(s)\to 2/A=e^{-\gamma}\), because \(a^\circ,c^\circ\)
are bounded by \(c_0\).
There can be no jump at \(2\) relative to the continuous Buchstab
solution.  Indeed, a jump \(\delta\) in \(sU(s)\) there would give
a difference \(\delta H(s)\), where \(H(s)=1/s\) on \([2,3]\) and
\((sH(s))'=H(s-1)\) thereafter.  The bound \(H\ge1/3\) on
\([2,3]\) persists by induction: if it holds up to an integer \(n\ge3\),
then for \(n\le s\le n+1\),
\[
 sH(s)=nH(n)+\int_n^s H(v-1)\,dv\ge n/3+(s-n)/3.
\]
For either sign of \(\delta\ne0\), this contradicts
\(U(s)-\omega_{\rm Bu}(s)\to0\), using the independently proved Buchstab limit in
Lemma~\ref{lem:elementary-buchstab}.
Thus \(U=\omega_{\rm Bu}\).  Together with
\(a^\circ+c^\circ=c_0\) and \eqref{eq:functions-gap-sum}, this identifies
\(a^\circ=a,\ c^\circ=c\), proving
\eqref{eq:functions-deficits} and both future integrals for \(s\ge2\).

Finally \(c(s)=A/s-1\) and \(a=1\) below \(2\).
The equation at \(s=2\) gives
\(\int_1^\infty a=A-2\); hence for \(1\le s\le2\),
\[
 \int_{s-1}^{\infty}a(t)\,dt=2-s+A-2=A-s=s c(s).
\]
This proves the remaining range of \eqref{eq:functions-future-c}.
\end{proof}
\subsection{Derivative weights and change of measure}

The deficits are now positive and have the required decay.  We use
derivatives of the sieve functions to normalize the harmonic transitions
of the standard tree.

Define
\[
 \phi_\ee(s)=s^2f'(s)\quad(s\ge2),\qquad
 \phi_\oo(s)=-s^2F'(s)\quad(s\ge0.95).
\]
At \(s=2\) use the right derivative, and on \(1.98\le s<2\) use
the derivative of \eqref{eq:functions-starting-extension}.
The state weight will be \(\phi_i(r/b)/r^2\); its reciprocal at the
endpoint is the factor that appears when we convert a probability
estimate back to harmonic weight.  For paths from the root, this gives
the scale \(B^{-2}\) when the endpoint gap and ratio remain in fixed
compact ranges.
Put
\[
 W_0(t)=\frac{t+1}{t^2},\qquad
 m_\ee(s)=s-1,\qquad m_\oo(s)=\max(2,s-1).
\]

\begin{lemma}[Derivative weights]\label{lem:derivative-weights}
The derivative weights are positive and nonincreasing on their indicated
ranges.  If \(i'\) denotes the type opposite to \(i\), then
\begin{equation}\label{eq:functions-derivative-integrals}
 \phi_i(s)=\int_{m_i(s)}^{\infty}W_0(t)\phi_{i'}(t)\,dt.
\end{equation}
This holds also on the specified starting extensions.
On a tail,
\begin{equation}\label{eq:functions-phi-dickman}
 \phi_\ee(s)\asymp\rho(s-2),\qquad
 \phi_\oo(s)\asymp\rho(s-2).
\end{equation}
There are absolute constants \(C,q\) such that the absolute logarithmic
derivatives, where they exist, and the ratios
\(\phi_{i'}(m_i(s))/\phi_i(s)\) are at most \(C(1+s)^q\).
In particular, for \(u,v\le S\) in the indicated domain of the
same weight, with \(S\ge3\),
\[
 |\log\phi_i(u)-\log\phi_i(v)|
 \le C(1+S)^q|u-v|.
\]
For each fixed \(\lambda>0\), the negative logarithmic derivatives
of both weights are at least \(\lambda\) on a sufficiently remote tail.
\end{lemma}

\begin{proof}
On the ordinary ranges the delay equations give
\begin{equation}\label{eq:functions-phi-deficits}
 \phi_\ee(s)=s\bigl(a(s)+c(s-1)\bigr),\qquad
 \phi_\oo(s)=s\bigl(c(s)+a(s-1)\bigr).
\end{equation}
The second expression uses \(a=1\) on \([0,1]\).
On the extensions, the explicit formulas are
\[
 \phi_\ee(s)=A\left(\frac{s}{s-1}-\log(s-1)\right),\qquad
 \phi_\oo(s)=A\quad(0.95\le s\le3).
\]
These formulas and Lemma~\ref{lem:future-integrals} prove positivity.
Differentiation gives
\begin{align}
 \phi_\ee'(s)&=-\frac{s}{(s-1)^2}\phi_\oo(s-1)
       &&(s\ge1.98),\label{eq:functions-phi-e-derivative}\\
 \phi_\oo'(s)&=-\frac{s}{(s-1)^2}\phi_\ee(s-1)
       &&(s>3),\label{eq:functions-phi-o-derivative}
\end{align}
at points of differentiability; the odd weight is constant below \(3\).
Equation \eqref{eq:functions-phi-deficits} and
\eqref{eq:functions-deficits} give \eqref{eq:functions-phi-dickman},
and in particular both weights vanish at infinity.
Integrating \eqref{eq:functions-phi-e-derivative} and
\eqref{eq:functions-phi-o-derivative} backwards from infinity gives
\eqref{eq:functions-derivative-integrals}.
For the even extension, the lower endpoint lies below \(1\), where
the explicitly constant odd weight gives exactly the displayed
continuation.  For odd \(s\le3\), the lower endpoint is \(2\), so the
integral is constant, as required.

On a tail, \eqref{eq:functions-dickman-ratio} and
\eqref{eq:functions-phi-dickman} give
\[
 \frac{\phi_{i'}(s-1)}{\phi_i(s)}=O(s^2).
\]
The derivative formulas then bound the absolute logarithmic
derivatives by a polynomial (in fact \(O(s)\) on a tail).
On the remaining compact ranges the explicit positive functions
and their piecewise derivatives give fixed bounds.
The weights are locally absolutely continuous by
\eqref{eq:functions-derivative-integrals}; integrating the
logarithmic derivative gives the stated Lipschitz bound.
Finally the same comparison gives a positive constant times
\[
 \frac{\rho(s-3)}{s\rho(s-2)}
 =\frac{s-2}{s}\left(-\frac{\rho'(s-2)}{\rho(s-2)}\right)
\]
as a lower bound for either negative logarithmic derivative on a
tail.  Apply \eqref{eq:functions-dickman-logderivative}.
\end{proof}

A \emph{standard} path replaces the admission condition
\eqref{eq:admission} by \(r-x\ge2x\) at an upper node.
Thus it contains every path in the original tree, and the two rules
agree for \(x\ge2\).
At a state of type \(i\), gap \(r\), cutoff \(b\), and ratio \(s=r/b\),
make a draw \(t\) with probability density
\begin{equation}\label{eq:tilted-kernel}
 K_i(s,dt)=
 \frac{W_0(t)\phi_{i'}(t)}{\phi_i(s)}
 \1_{\{t\ge m_i(s)\}}\,dt.
\end{equation}
Lemma~\ref{lem:derivative-weights} proves that this is a probability
kernel.  The next exponent, gap, and ratio are
\[
 x=\frac r{t+1},\qquad r'=\frac{rt}{t+1},\qquad s'=t.
\]
The lower support endpoint ensures \(x\le b\), and the upper-to-lower
transition also satisfies \(s'\ge2\).
Put \(g(t)=\log(1+1/t)\), so \(r'=r e^{-g(t)}\).

The harmonic transition \(dx/x\), after reversing orientation in
\(t=r/x-1\), is \(dt/(t+1)\).  Dividing this by
\eqref{eq:tilted-kernel} gives
\[
 \frac{t^2}{(t+1)^2}\frac{\phi_i(s)}{\phi_{i'}(t)}
 =\left(\frac{r'}r\right)^2
    \frac{\phi_i(s)}{\phi_{i'}(t)}.
\]
Consequently, for every finite prefix,
\begin{equation}\label{eq:path-weight}
 \text{original harmonic weight}
 =
 \text{probability weight}\,
 \left(\frac{r_{\rm end}}{r_{\rm in}}\right)^2
 \frac{\phi_{i_{\rm in}}(s_{\rm in})}
      {\phi_{i_{\rm end}}(s_{\rm end})}.
\end{equation}
For a nonempty prefix its last exponent is
\(x=r_{\rm end}/s_{\rm end}\).  Since successive cutoffs are
nonincreasing, restriction of all exponents to \(x>\ell\) is
equivalent to the single arrival condition
\begin{equation}\label{eq:functions-arrival-cutoff}
 s_{\rm end}<r_{\rm end}/\ell.
\end{equation}

In the rest of this section, starts are of even type with
\begin{equation}\label{eq:functions-start-range}
 s_{\rm in}\in[1.99,2.3].
\end{equation}
Let \((i_n,s_n)\) be the resulting infinite sequence, \(n=0,1,\ldots\),
and put
\[
 T_0=0,\qquad T_n=\sum_{j=1}^n g(s_j),\qquad
 r_n=r_{\rm in}e^{-T_n},\qquad x_n=r_n/s_n\quad(n\ge1).
\]
After two steps the state spaces are
\(\{\ee\}\times[2,\infty)\) and \(\{\oo\}\times[1,\infty)\).
Only the first odd state can lie below \(1\).

\subsection{Regeneration and invariant occupation}

Call an odd state with \(s\le3\) a regenerating state; regeneration is
viewed as occurring just before its next draw.
Its outgoing law is independent of \(s\), since
\(\phi_\oo(s)=A\) and \(m_\oo(s)=2\).

\begin{lemma}[Regeneration]\label{lem:regeneration}
The time to first regeneration has an exponential moment uniformly
over \eqref{eq:functions-start-range}.  The cycles strictly after one
regeneration through the next inclusive are iid.  Their lengths
\(\tau\) have a finite exponential moment.
Their total log decrements
\[
 G=\sum_{j=1}^{\tau}g(s_j)
\]
are positive, nonarithmetic, and have a finite exponential moment for
some positive parameter.  In particular \(0<\E G<\infty\).
Each complete cycle satisfies \(G\ge g(3)=\log(4/3)\).
\end{lemma}

\begin{proof}
For sufficiently large current \(s\), the excess of the next state
over \(s-1\) has tail
\[
 \Pbb(s_{\rm next}\ge s-1+h\mid i,s)
 =\frac{\phi_i(s+h)}{\phi_i(s)}\qquad(h\ge0).
\]
Here the support minimum is \(s-1\) for either type, and the identity
follows from \eqref{eq:functions-derivative-integrals}.
Lemma~\ref{lem:derivative-weights} makes this tail at most
\(e^{-\lambda h}\), with any fixed large \(\lambda\), outside a
sufficiently large compact.  Thus \(V(i,s)=e^s\) satisfies
\[
 \E(V(i_{\rm next},s_{\rm next})\mid i,s)
 \le \kappa V(i,s)+C,\qquad \kappa<1,
\]
and the constant term can be omitted outside a larger compact.
On the compact, finiteness and uniform boundedness of the expectation
follow from the superexponential tails of the weights.

Choose a compact set \(C_0\) and \(\alpha>1\) such that
\(\E(V_{\rm next}\mid i,s)\le V(i,s)/\alpha\) outside \(C_0\).
Stopping the nonnegative supermartingale
\(\alpha^{n\wedge\sigma}V(i_{n\wedge\sigma},s_{n\wedge\sigma})\),
where \(\sigma\) is the first entrance to \(C_0\), gives
\[
 \E_{i,s}\alpha^\sigma\le V(i,s)
\]
for starts outside \(C_0\), by Fatou's lemma.
In particular the entrance time is almost surely finite.

From \(C_0\) there are an integer \(m\) and a probability \(p>0\)
such that regeneration is reached within \(m\) steps with probability
at least \(p\), uniformly in the starting state.
To see this, whenever \(s>3\) choose the next state in the interval
of width \(1/4\) immediately above its support minimum; this lowers
the state by at least \(3/4\).
An even state with \(s\le3\) then draws into the regenerating odd
range on the same choice.
All prescribed intervals lie in a fixed compact and have uniformly
positive transition probabilities.

Run such an \(m\)-step attempt, and on failure wait for the next
entrance to \(C_0\) before retrying.  Let \(D\) be the duration of
one failed attempt together with this return.  The stopped estimate
and the uniform bound on the \(V\)-moment after the \(m\) attempted
steps show that, for some \(\alpha_1>1\),
\[
 \sup_{(i,s)\in C_0}\E_{i,s}
       [\alpha_1^D\1_{\{\text{failure}\}}]<\infty.
\]
For \(0<\theta<1\), Hölder's inequality bounds the same expression
with \(\alpha_1^\theta\) by
\[
 \left(\sup_{C_0}\E[
       \alpha_1^D\1_{\{\text{failure}\}}]\right)^\theta
 (1-p)^{1-\theta}.
\]
For sufficiently small \(\theta>0\) this is less than \(1\).
Summing over successive failed attempts proves an exponential
moment for the hitting time of regeneration from \(C_0\), and then
from \eqref{eq:functions-start-range}.
The same reasoning applies to a return after a regenerating start:
its outgoing common draw has a bounded \(V\)-moment.

The common outgoing law proves independence and identical
distribution of the complete cycles.  During a cycle all odd
states are at least \(1\), and all even states are at least \(2\),
so each \(g(s_j)\le\log2\).
Thus the exponential moment of \(\tau\) supplies an exponential
moment of \(G\) for a sufficiently small parameter.
The terminal odd state lies in \([1,3]\), so its last decrement
alone is at least \(g(3)>0\).

For nonarithmeticity, consider a two-step return whose even state
\(t\) lies in a fixed subinterval of \((2,3)\) and whose next odd
state \(u\) lies strictly between \(t-1\) and \(3\).
There is a two-dimensional region of such pairs on which the
joint density is positive.  On that region
\(G=g(t)+g(u)\) varies continuously and has a nonzero derivative
in \(u\).  Its distribution therefore has a nonzero absolutely
continuous component, and cannot be supported on an arithmetic
lattice.
\end{proof}

\begin{lemma}[Invariant densities and cycle occupation]
\label{lem:invariant}
On the ordinary even and odd state spaces, respectively, the finite
raw invariant densities are
\begin{equation}\label{eq:functions-invariant-densities}
 \pi_\ee(s)=\phi_\ee(s)\quad(s\ge2),\qquad
 \pi_\oo(s)=(1-s^{-2})\phi_\oo(s)\quad(s\ge1).
\end{equation}
Put
\[
 \beta=\int_1^3\pi_\oo(s)\,ds=\frac{4A}{3}.
\]
For every nonnegative measurable state function \(Q\),
\begin{equation}\label{eq:functions-cycle-occupation}
 \E\sum_{j=1}^{\tau}Q(i_j,s_j)
 =\frac1\beta\sum_i\int Q(i,s)\pi_i(s)\,ds.
\end{equation}
In particular, if
\[
 M_g=\sum_i\int\pi_i(s)g(s)\,ds,
\]
then \(\E G=M_g/\beta\), and
\begin{equation}\label{eq:mg-bound}
 0<M_g\le4e^\gamma.
\end{equation}
\end{lemma}

\begin{proof}
For a new even state \(t\ge2\), cancellation of the old derivative
weight in the incoming integral gives
\[
 \int_1^{t+1}(1-s^{-2})\,ds=\frac{t^2}{t+1}.
\]
Multiplication by \(W_0(t)\phi_\ee(t)\) yields \(\pi_\ee(t)\).
For a new odd state \(t\ge1\), the corresponding integral is
\(\int_2^{t+1}ds=t-1\), which yields
\((1-t^{-2})\phi_\oo(t)\).  These identities prove invariance.
Finiteness follows from \eqref{eq:functions-phi-dickman}.
On \([1,3]\), \(\phi_\oo=A\), so
\(\beta=A[s+s^{-1}]_1^3=4A/3\).

For completeness, normalize the raw measure to a stationary
probability measure \(\mu=\pi/M\), where
\(M=\sum_i\int\pi_i\).
Start the chain with law \(\mu\).
For a bounded measurable state function \(Q\), partition \(Q(i_n,s_n)\)
according to the last regenerating time strictly before \(n\).
If it is \(n-k\), \(1\le k\le n\), the common outgoing law shows
that this term has expectation
\[
 \mu(\text{regenerating states})\,
 \E[Q(i_k,s_k)\1_{\{\tau\ge k\}}],
\]
where the expectation on the right starts just after a regeneration.
The remaining event is that there was no regenerating time among
\(0,\ldots,n-1\); its probability tends to zero by
Lemma~\ref{lem:regeneration} and dominated convergence over \(\mu\).
Stationarity and then \(n\to\infty\) give
\[
 \mu(Q)=\frac{\beta}{M}
       \E\sum_{j=1}^{\tau}Q(i_j,s_j).
\]
This proves \eqref{eq:functions-cycle-occupation} for bounded
nonnegative \(Q\), and truncation proves it in general.
The argument does not require an aperiodic discrete-time chain;
the alternating types cause no difficulty.

Since \(g(s)\le1/s\), \eqref{eq:functions-phi-deficits} gives
\[
 M_g\le
 \int_2^\infty a(s)\,ds+
 2\int_1^\infty c(s)\,ds+
 \int_0^\infty a(s)\,ds.
\]
The future integrals at \(s=2\) give
\(\int_1^\infty c=2\) and \(\int_1^\infty a=A-2\).
As \(a=1\) on \([0,2]\), the three displayed contributions are
\(A-3,\ 4,\ A-1\).  Their sum is \(2A=4e^\gamma\).
Strict positivity is immediate from the positive even density and
\(g>0\).  Taking \(Q(i,s)=g(s)\) in
\eqref{eq:functions-cycle-occupation} identifies \(\E G\).
\end{proof}

\subsection{Renewal limits and uniform band estimates}

We use tests in the log-gap coordinate \(u=\log r\).
An admissible compact test \(F_i(u,s)\) is bounded and continuous on
the extended state ranges, with support in a fixed compact rectangle
in \((u,s)\).  We also permit multiplication by indicators of fixed
intervals in \(u\) or \(s\), and by
\(\1_{\{s<e^u/\ell\}}\), where \(\ell=1\) or \(2\).
At the lower state endpoints one may first make a continuous
one-sided extension to the extended range and then apply the
indicator.  Finite linear combinations of these tests are permitted.
In particular, a continuous bounded test supported in a compact
positive \(r\)-window and restricted by \(s<r/\ell\) is covered:
the restriction itself bounds \(s\).

\begin{proposition}[Compact occupation limit]\label{prop:occupation}
For an admissible compact test and starts satisfying
\eqref{eq:functions-start-range},
\begin{equation}\label{eq:functions-occupation-limit}
 \lim_{v\to\infty}
 \E_{s_{\rm in}}\sum_{n\ge0}F_{i_n}(v-T_n,s_n)
 =
 \frac1{M_g}\sum_i\int_{\R}\int F_i(u,s)\pi_i(s)\,ds\,du.
\end{equation}
The convergence is uniform in \(s_{\rm in}\in[1.99,2.3]\).
The state integrals on the right use the ordinary invariant ranges
in \eqref{eq:functions-invariant-densities}.
\end{proposition}

\begin{proof}
First start immediately before the outgoing draw of a regeneration,
and exclude that starting state from the sum.
Write a complete cycle as
\((i_j,s_j)_{1\le j\le\tau}\), and set
\(S_j=\sum_{k=1}^j g(s_k)\) within the cycle.
Its response to the test is
\begin{equation}\label{eq:functions-cycle-test}
 H_0(v)=\E\sum_{j=1}^{\tau}F_{i_j}(v-S_j,s_j).
\end{equation}
If the log-gap support of the test is contained in \([a,b]\),
then \(H_0(v)=0\) for \(v<a\), since \(S_j\ge0\).
Also \(S_j\le j\log2\), so for \(v>b\),
\[
 |H_0(v)|\le
 \|F\|_\infty
 \E\!\left[\tau\,
 \1_{\{\tau\ge(v-b)/\log2\}}\right].
\]
The exponential moment in Lemma~\ref{lem:regeneration} gives an
exponentially decreasing upper bound as \(v\to\infty\).

The function \(H_0\) is continuous, including for the permitted
indicator tests.  To verify this, condition on the history before
the \(j\)-th draw and on the event \(j\le\tau\), which is determined
before that draw.  The next state \(t=s_j\) has a density.
Writing \(D=e^{v-S_{j-1}}>0\), its gap is
\[
 e^{v-S_j}=\frac{Dt}{t+1}.
\]
Equality at a fixed positive gap edge therefore specifies at most
one value of \(t\); equality at a fixed state edge does likewise.
For the arrival boundary,
\[
 t=e^{v-S_j}/\ell
 \quad\Longleftrightarrow\quad
 t+1=D/\ell.
\]
Thus each boundary has conditional probability zero.
For any fixed number of cycle steps dominated convergence gives
continuity.  The remaining steps are uniformly bounded in expectation
by \(\|F\|_\infty\E[\tau\1_{\{\tau>N\}}]\), which tends to zero.
This proves continuity of \(H_0\).
Continuity, the vanishing left tail, and the exponential right
envelope show that \(H_0\) is directly Riemann integrable.

Let \(U_G\) be the renewal measure of successive complete-cycle
decrements, including the boundary at zero:
\[
 U_G=\sum_{k\ge0}\mathcal L(G_1+\cdots+G_k).
\]
The expectation of the complete future test from a regeneration is
\begin{equation}\label{eq:functions-renewal-convolution}
 R_F(v)=\int_{[0,\infty)}H_0(v-y)\,U_G(dy).
\end{equation}
By Lemma~\ref{lem:regeneration}, \(G\) is nonarithmetic and has
finite positive mean.  The key renewal theorem
\eqref{eq:renewal-input} therefore gives
\[
 R_F(v)\longrightarrow \frac1{\E G}\int_\R H_0(u)\,du.
\]
Fubini is justified by the bounded test, its compact log-gap
support, and \(\E\tau<\infty\).  Translating the log-gap integral
within each cycle and applying
\eqref{eq:functions-cycle-occupation} gives
\[
 \int_\R H_0(u)\,du
 =\E\sum_{j=1}^{\tau}\int_\R F_{i_j}(u,s_j)\,du
 =\frac1\beta\sum_i\int\!\!\int F_i(u,s)\pi_i(s)\,ds\,du.
\]
Since \(\E G=M_g/\beta\), this is the limit in
\eqref{eq:functions-occupation-limit}.

We record a boundedness fact needed to restore the initial segment.
Every complete cycle has decrement at least \(g(3)>0\).
Consequently any interval of fixed length contains a bounded number
of cycle boundaries on each sample path, and
\begin{equation}\label{eq:functions-renewal-band}
 \sup_{y\in\R}U_G([y,y+h])\le
 1+\left\lfloor\frac{h}{g(3)}\right\rfloor
 \qquad(h>0).
\end{equation}
The compact bound and the exponential envelope for \(H_0\), summed
over unit intervals, show from
\eqref{eq:functions-renewal-convolution} that
\(\sup_v|R_F(v)|<\infty\).

Now let \(\tau_0\) be the first regenerating time from an allowed
even start, and \(D_0=T_{\tau_0}\).
All initial per-step decrements are bounded by
\(c_*=\log(1+1/0.98)\).  Lemma~\ref{lem:regeneration} gives a
uniform exponential moment for \(\tau_0\), hence a uniformly
tight, indeed exponentially bounded, family of \(D_0\)'s.
The portion of the test through time \(\tau_0\) tends to zero
uniformly as \(v\to\infty\): a contribution requires
\(c_*\tau_0\ge v-b\), and its absolute size is at most
\(\|F\|_\infty(\tau_0+1)\).
Conditional on the initial history, the future after that state
has expectation \(R_F(v-D_0)\).
It is the same function \(R_F\) for every history, because the
outgoing regenerating law is common.
For a fixed \(M\), convergence of \(R_F(y)\) as \(y\to\infty\)
is uniform over \(y=v-d,\ 0\le d\le M\).
Uniform tightness of \(D_0\) and boundedness of \(R_F\) then give
the claimed uniform limit.
\end{proof}

\begin{lemma}[Uniform band bounds]\label{lem:band-bounds}
For every fixed \(h>0\) there is \(C_h<\infty\) with the following
properties.  For any nonnegative measurable state function \(Q\)
and any \(v\in\R\), a regenerating start satisfies
\begin{equation}\label{eq:functions-band-from-regeneration}
 \E\sum_{n\ge1}Q(i_n,s_n)
        \1_{\{T_n\in[v,v+h]\}}
 \le C_h\sum_i\int Q(i,s)\pi_i(s)\,ds.
\end{equation}
From an even start in \eqref{eq:functions-start-range}, the same
bound holds for the sum over \(n\ge2\), uniformly in the starting
ratio.  The density of the first odd state is at most
\[
 C\,\1_{\{s\ge0.98\}}\phi_\oo(s).
\]
In particular, uniformly in \(r_{\rm in}>0\), \(R>0\), and
\(\ell\in\{1,2\}\),
\begin{equation}\label{eq:functions-inverse-phi-band}
 \E\sum_{n\ge1}
 \frac{\1_{\{R\le r_n\le eR,\ s_n<r_n/\ell\}}}
      {\phi_{i_n}(s_n)}
 \ll R+1.
\end{equation}
Including the starting state in the last sum changes its bound
by at most an absolute constant.
\end{lemma}

\begin{proof}
For a regenerating start, condition on the internal states and
decrements of one fresh cycle.  Its beginning is a renewal boundary
independent of that cycle.  Tonelli's theorem gives
\[
 \E\sum_{n\ge1}Q(i_n,s_n)\1_{\{T_n\in[v,v+h]\}}
 =
 \E\sum_{j=1}^{\tau}Q(i_j,s_j)
 U_G([v-S_j,v+h-S_j]).
\]
Apply \eqref{eq:functions-renewal-band} and then
\eqref{eq:functions-cycle-occupation}.
This argument permits arbitrary nonnegative \(Q\): it never
assumes independence of a state's value and its internal decrement.

We next compare the initial even start to regeneration.
If its ratio is \(s_0\), the density after two draws, on \(t\ge2\),
is
\begin{equation}\label{eq:functions-two-step-density}
 p_{s_0,2}(t)=
 \frac{W_0(t)\phi_\ee(t)}{\phi_\ee(s_0)}
 \int_{s_0-1}^{t+1}W_0(u)\,du.
\end{equation}
Starting at regeneration, the density after three draws is
\begin{equation}\label{eq:functions-three-step-density}
 p_{{\rm reg},3}(t)=
 \frac{W_0(t)\phi_\ee(t)}A
 \int_1^{t+1}W_0(u)
       \int_2^{u+1}W_0(y)\,dy\,du
 \qquad(t\ge2).
\end{equation}
These formulas follow by canceling the intermediate derivative
weights in the successive kernels.
For \(u\ge2\), the inner integral in
\eqref{eq:functions-three-step-density} is at least
\(\int_2^3W_0(y)\,dy>0\).
Since \(t+1\ge3\), it follows that
\[
 \int_1^{t+1}W_0(u)\int_2^{u+1}W_0(y)\,dy\,du
 \gg \int_2^{t+1}W_0(u)\,du
 \gg \int_{s_0-1}^{t+1}W_0(u)\,du,
\]
uniformly for \(s_0\in[1.99,2.3]\).
The last comparison uses the bounded integral over
\([0.99,2]\) and the positive lower bound for
\(\int_2^{t+1}W_0\).
The starting denominator \(\phi_\ee(s_0)\) also has a positive
compact lower bound.  Thus
\begin{equation}\label{eq:functions-initial-density-domination}
 p_{s_0,2}(t)\le C p_{{\rm reg},3}(t)\qquad(t\ge2).
\end{equation}

To use this marginal comparison, the correlated elapsed decrement
must be removed explicitly.  Let \(I=[v,v+h]\),
\(C_2=2c_*\), and \(C_3=3\log2\).
The initial two decrements sum to at most \(C_2\), so an event
in \(I\) at time \(n\ge2\) has, measured from state \(2\), a future
decrement in
\[
 I_2=[v-C_2,v+h].
\]
Condition on that state, discard the initial decrement by this
enlargement, and apply
\eqref{eq:functions-initial-density-domination}.
For a path started at regeneration, its first three decrements
sum to at most \(C_3\); a future decrement in \(I_2\) measured
from state \(3\) therefore gives a total decrement in
\[
 [v-C_2,v+h+C_3].
\]
The resulting expectation is bounded by a constant times the
regenerating occupation sum in this enlarged band, which is
covered by \eqref{eq:functions-band-from-regeneration}.
This proves the bound from even starts, including their state
at time \(2\), without imposing independence on the initial
state and elapsed decrement.

At step \(1\), the density is
\[
 \frac{W_0(s)\phi_\oo(s)}{\phi_\ee(s_0)}
 \1_{\{s\ge s_0-1\}}.
\]
Both \(W_0(s)\) for \(s\ge0.99\) and
\(1/\phi_\ee(s_0)\) on the allowed starts are bounded.
This proves the asserted first-step estimate.

The condition \(R\le r_n\le eR\) is a band of length \(1\) in
\(T_n\).  For \(n\ge2\), use the preceding result with
\[
 Q(i,s)=\frac{\1_{\{s<eR/\ell\}}}{\phi_i(s)}.
\]
The ratios of invariant to derivative densities are \(1\) for
the even states and \(1-s^{-2}\le1\) for the odd states.
Their integrals up to \(eR/\ell\) are \(O(R+1)\).
The first-step density gives the same bound by direct integration.
Finally \(\phi_\ee(s_0)\) has a positive compact lower bound, so
the starting state's possible contribution is \(O(1)\).
\end{proof}

\begin{remark}\label{rem:functions-compact-weight}
Two consequences explain the later weight scales.
At an arrival with \(x_n>\ell\) and \(r_n\le K\),
\eqref{eq:functions-arrival-cutoff} confines the state to a
fixed compact with a positive lower endpoint.  Hence
\(r_n^2/\phi_{i_n}(s_n)\) is bounded.
Moreover, each such transition decreases the gap by more than
\(\ell\); there are only \(O_K(1)\) such compact arrivals on
one path.  Formula \eqref{eq:path-weight} therefore bounds the
total original harmonic weight of compact prefixes by
\(O_K(r_{\rm in}^{-2})\).
If the counted prefixes can occur only on an event \(\mathcal A\)
of the tilted path, the bound is
\(O_K(r_{\rm in}^{-2})\Pbb(\mathcal A)\).

Retain the restriction \(x_n>\ell\), equivalently
\(s_n<r_n/\ell\), in the following tail estimate.
After multiplication by
\(r_{\rm in}^2/\phi_\ee(s_{\rm in})\), a reward \(e^{-c r}\)
in original harmonic weight is \(r^2e^{-c r}/\phi_i(s)\)
in tilted weight.  Summing
\eqref{eq:functions-inverse-phi-band} over the bands
\([R,eR]\), \(R\ge K\), gives a convergent upper bound whose
tail tends to zero as \(K\to\infty\), uniformly in the allowed
starts.  These statements also cover a fixed bounded factor
in the reward.
\end{remark}

\subsection{Visits in a prescribed even-state window}

\begin{lemma}[Marked visits]\label{lem:marked-visits}
Fix \(b_0>0\) and \(\eps>0\).  There is a constant
\(C_\eps>10\), independent of \(b_0\), such that, whenever
\(b_1\ge C_\eps b_0\) is fixed, the following event has
probability at least \(1-\eps-o(1)\) as \(r_{\rm in}\to\infty\),
uniformly in \eqref{eq:functions-start-range}.
At a regenerating odd state with pre-draw gap in
\([10b_0,b_1]\), the next draw is an even arrival satisfying
\[
 s\in[2.08,2.14],\qquad x\in[2b_0,b_1/2].
\]
The proof yields the smaller ratio interval \([2.09,2.13]\);
thus the stated ratio window has positive slack.
The arriving gap exceeds \(6b_0\), so, if \(6b_0>K\), such a visit
precedes every subsequent arrival of gap at most \(K\).
\end{lemma}

\begin{proof}
Mark a complete regeneration cycle if its first draw lies in
\([2.09,2.13]\).  The common outgoing kernel gives a fixed
probability
\[
 p=\frac1A\int_{2.09}^{2.13}W_0(t)\phi_\ee(t)\,dt
 \in(0,1).
\]
The complete cycle objects, consisting of their state sequences
and their decrements, are iid.
Consequently the spacings \(W\) between successive marked cycle
beginnings are iid: each consists of one cycle conditioned on
being marked, followed by a geometric number of cycles conditioned
on being unmarked.  More explicitly, the number \(N\) of intervening
unmarked cycles has
\(\Pbb(N=m)=(1-p)^m p\), \(m\ge0\); conditional on these marks,
the cycle objects are independent with their respective
conditional laws.
The next marked cycle is excluded from the preceding spacing
and begins the next such block.  This proves independence of
the spacings, even though a mark depends on its own cycle.
It also gives
\[
 \E W=
 \E(G\mid\text{marked})
 +\frac{1-p}{p}\E(G\mid\text{unmarked})
 =\frac{\E G}{p}<\infty.
\]
Furthermore \(W\ge g(3)>0\).

From an allowed initial start, the delay \(D\) to the first
marked beginning consists of the initial segment to regeneration
and a geometric number of unmarked complete cycles.
The first segment has uniformly bounded mean by
Lemma~\ref{lem:regeneration}; the latter sum has finite mean by
the preceding conditional-cycle description.
Thus \(D\) is uniformly tight over the allowed initial states.
The spacings after that beginning are independent of the delay.

We give the required wide-interval bound using only
\(\E W<\infty\).
Let \(U_W\) be the renewal measure for the marked spacings,
with an initial renewal at zero.
The lower bound on \(W\) gives a uniform bound \(C\) for its
mass in any interval of length one.
For \(q\ge0\), if there is no marked renewal in \([q,q+H]\),
some renewal at or before \(q\) has its next spacing overshoot
\(q+H\).  Independence of a renewal time and its next spacing
therefore gives
\[
 \Pbb(\text{no renewal in }[q,q+H])
 \le\int_{[0,q]}\Pbb(W>q+H-t)\,U_W(dt)
 \le C\sum_{j\ge0}\Pbb(W>H+j).
\]
Changing one endpoint in the unit-interval partition only changes
the harmless absolute constant or shifts \(H\) by \(1\).
The last quantity tends to zero as \(H\to\infty\), since \(W\)
has finite mean.  It is uniform in \(q\).
For the delayed process, condition on \(D\le v\) and apply the
same bound with \(q=v-D\).  The additional error
\(\Pbb(D>v)\) tends uniformly to zero as \(v\to\infty\).

Apply this with
\[
 v=\log(r_{\rm in}/b_1),\qquad
 H=\log\bigl(b_1/(10b_0)\bigr).
\]
Choose \(H\) sufficiently large in terms of \(\eps\), equivalently
choose \(b_1/b_0\) sufficiently large, and then let
\(r_{\rm in}\to\infty\).
A marked beginning in this log-decrement interval has pre-draw
gap \(r\in[10b_0,b_1]\).
For its marked draw \(t\in[2.09,2.13]\),
\[
 \frac{10b_0}{3.13}\le x=\frac r{t+1}
 \le\frac{b_1}{3.09}.
\]
This lies inside \([2b_0,b_1/2]\), and the arriving ratio is \(t\).
The new gap satisfies
\[
 r'=\frac{rt}{t+1}\ge
 10b_0\frac{2.09}{3.09}>6b_0.
\]
Gaps strictly decrease along the path, so the last assertion
follows.
\end{proof}
